\begin{proof}
Here $3/5\le a<r$. The inward successors lie in $[1-a,a]$;
the two outward ones are clipped to their respective endpoints.
Their barycentric coefficients give the displayed matrices.
The aggregate has both row and column sums $1/2$, hence $x_L=(1,1)$.
If $f^{[1]},f^{[2]}$ are the laws from the two original pure states,
then $(1,-1)A_0=(1/10,0)$ and $(1,-1)A_1=(0,-1/10)$.
Their difference has positive mass $1/10$ on words starting with $0$
and negative mass $-1/10$ on words starting with $1$, with zero
empty-word difference. Thus
$\|f_t-f_s\|_1=|t-s|/5$.
The outward coordinate displacement is
$(1-a-a^2)/(2+a)$ and its letter probability is $(2+a)/10$.
The two occupied rows therefore give $(1-a-a^2)/25$.
The cross-ratio formula gives
$r-a=\sqrt5\,z_L/(1+z_L)$, so
$1-a-a^2=5z_L/(1+z_L)^2$.
Theorem~\ref{II:U:thm:flows} completes the proof.
\end{proof}
